\section{Design}
\label{sec:design}

\sys{} is an OS-level harness for agent policy compliance. Its central
design choice is to let agents and project owners express policy in terms
of task and repository intent, while enforcing the resulting constraints
below the tool layer. This separates the context needed to instantiate a
policy from the mechanism that observes system-level behavior.

The design challenge is to keep these two sides connected without
collapsing one into the other: a policy must be concrete enough for
deterministic OS-level enforcement, but expressive enough to encode
project and task intent supplied by the agent. It must also preserve
cross-event state and prevent lower-authority policy from weakening
higher-authority constraints.

\subsection{Threat Model and Assumptions}

\sys{} targets a cooperative but unreliable agent. The agent is expected
to declare task intent and project constraints honestly, but may fail to
follow them because of planning errors, prompt injection, shell-script
side effects, or opaque tool behavior. The system does not target a
malicious process that attacks the kernel, exploits the loader, or
deliberately evades policy through adversarial obfuscation.

The trusted computing base consists of the compiler, the eBPF enforcement
engine, the runner, and higher-authority policies. The agent's process
tree, including tools, shells, subprocesses, and sub-agents, is treated as
the monitored domain. Out of scope are kernel compromise, root or
\texttt{CAP\_BPF} compromise, deliberate side channels, and semantic
content checks that are not observable at process, file, or network
edges.

\subsection{System Overview}

The empirical study in \S\ref{sec:empirical} gives three design
requirements. First, policy must be \emph{agent-facing}: many directives
refer to project- or task-specific concepts that only the agent can
resolve in context. Second, enforcement must be \emph{system-level}:
directives often constrain commands, files, subprocesses, and network
connections rather than final model text. Third, enforcement must be
\emph{stateful}: many directives depend on previous actions, such as
whether tests ran after a source change or whether data from one source
later reached another sink.

\sys{} implements these requirements with three abstractions. A compact
DSL lets agents and maintainers bind natural-language directives to
source, sink, and condition patterns. A labeled information-flow model
maintains cross-event state across processes, files, and network
endpoints. A layered authority model lets lower-authority policy add or
narrow constraints without weakening platform or project invariants.

\subsection{Policy DSL}

The DSL is the boundary between intent-level instructions and
system-level enforcement. A policy contains sources that introduce labels,
rules that match labeled events at sinks, and optional gates that express
ordering or controlled declassification. This keeps the agent-facing
surface small: translating an instruction usually requires producing one
declarative rule rather than constructing the backend representation used
by the kernel.

Sources bind labels to system objects. An \texttt{exec} source labels
processes that execute a matching binary, a \texttt{file} source labels
matching paths, and a network source labels matching endpoints. Rules bind
an effect such as \texttt{notify}, \texttt{block}, or \texttt{kill} to an
operation such as \texttt{exec}, \texttt{read}, \texttt{write},
\texttt{unlink}, or \texttt{connect}. Conditions express common agent
workflow constraints: a target allow-list, a lineage gate, or a temporal
\texttt{after ... since ...} gate.

\begin{figure}[t]
\begin{verbatim}
source AGENT = exec "claude"
source SECRET = file "**/.env"

# Per-event: from CoplayDev/unity-mcp
rule no-push-main:
  kill exec "git" "push" "main" if AGENT
  because "Do not push to main; open a PR instead."

# Cross-event: from chenhg5/cc-connect
rule tests-before-commit:
  block exec "git" "commit"
    if AGENT unless after exec "go" since write "**/*.go"
  because "Run `go test ./...` before committing."
\end{verbatim}
\caption{Two \sys{} rules drawn from real projects. The first is a
per-event rule; the second is a cross-event rule that requires temporal
state.}
\Description{Code listing of two \sys{} rules from real projects.}
\label{fig:policy-example}
\end{figure}

Figure~\ref{fig:policy-example} shows the two main rule shapes. The first
case is local to one event: if the agent process tries to push to
\texttt{main}, the event matches immediately. The second case depends on
history: committing is compliant only if the required test command ran
after the relevant source write. This distinction is why the DSL exposes
both per-event predicates and cross-event gates.

The DSL is intentionally narrower than a general mandatory access-control
language. It focuses on policy shapes that recur in agent instruction
files: do not use a command, do not write outside a scope, do not connect
to a sink after reading a source, run a validation command before a
release action, and route sensitive data through an approved tool. The
compiler lowers these rules into fixed-size kernel tables; policy authors
do not need to reason about BPF maps, verifier constraints, or binary
layout.

\subsection{Labeled Information-Flow Model}

\sys{} represents policy state as labels on OS objects. Processes, files,
and network endpoints each carry a 64-bit label mask. A source declaration
adds labels when an object matches the source pattern. Let $L(x)$ denote
the label mask of object $x$. Labels propagate through observed
system-level edges:

\begin{itemize}
\item \textbf{fork}$(p \to c)$: $L(c) := L(c) \cup L(p)$.
\item \textbf{exec}$(p, f)$: $L(p) := L(p) \cup L(f) \cup S(f)$, where
$S(f)$ is the set of source labels matching $f$.
\item \textbf{read}$(p, f)$: $L(p) := L(p) \cup L(f)$.
\item \textbf{write}$(p, f)$: $L(f) := L(f) \cup L(p)$.
\item \textbf{connect}$(p, e)$: $L(e) := L(e) \cup L(p)$.
\end{itemize}

These transfer functions maintain the key invariant: if an object carries
label $\ell$, then information from a source tagged with $\ell$ has
reached that object through observed process, file, or network edges.
Rules then check whether an event's subject or target contains required
labels, contains forbidden labels, or satisfies a temporal gate.

The model operates at object granularity rather than byte granularity,
trading precision for low-overhead whole-session coverage. This is a
deliberate fit for agent harnesses. Repository instructions usually refer
to objects such as commands, files, directories, repositories, endpoints,
and workflow steps, not individual bytes. Object-level labels can express
cross-event compliance patterns that tool-call filters miss, while
avoiding the complexity of dynamic byte-level taint
tracking~\cite{kemerlis2012libdft}.

Labels are monotonic by default: propagation adds labels and does not
remove them. When a project requires controlled release, the DSL can name
an explicit gate, such as a redaction command or review tool, that clears
specified labels for subsequent events. This makes declassification a
declared policy action rather than an implicit side effect of the agent's
plan.

\subsection{Layered Policy Authority}

Agent-authored policy is useful only if it cannot weaken stronger
constraints. \sys{} therefore treats policy as a layered authority stack:
platform/operator policy, project-maintainer policy, and agent/session
policy. Lower layers may add rules, specialize targets, add labels, or
narrow scope. They may not delete, rewrite, declassify, or weaken rules
introduced by a higher-authority layer.

\sys{} separates a pure IFC rule from the runtime binding that activates
it. A rule defines a condition, an effect, and a reason. A domain is the
policy boundary for a process tree; each process belongs to one domain,
and each event is checked against that domain's effective policy. Files
and endpoints are IFC objects that carry labels, not policy domains.

Bindings attach rules to domains as either \emph{locked} or
\emph{default}. Locked bindings are inherited security invariants and
cannot be disabled by child domains. Default bindings are templates that a
child domain may keep, disable for itself, or replace with stricter local
policy. This gives the agent room to adapt policy to a task while
preserving the authority of platform and project constraints.

The prototype resolves the active authority stack when policy is compiled
or loaded. The kernel receives the resulting rule tables and evaluates
events against them; it does not parse YAML, resolve project context, or
interpret authority metadata on the event path. This preserves the design
separation: agents and maintainers supply context and policy structure,
while the OS-level engine enforces the compiled result consistently across
tools and subprocesses.

\section{Implementation}
\label{sec:implementation}

\sys{} is implemented as a Rust compiler/runner and an eBPF enforcement
engine. The compiler parses the DSL, allocates label bits, lowers Boolean
label expressions and path/network patterns, and emits a fixed-size
C-compatible configuration consumed by the eBPF loader. Human-readable
rule names and reasons remain in userspace metadata; the kernel receives
compact rule tables and emits rule identifiers when events match.

The eBPF engine attaches to process, file, and network hooks. It maintains
label state for processes, files, and endpoints in BPF maps, applies the
transfer functions above, and checks the compiled rules on each observed
event. BPF-LSM hooks provide pre-operation enforcement for blocking
actions; tracepoints provide observation and post-operation termination
where LSM support is unavailable. Userspace does not re-detect policy
violations: it loads policy, starts the target process, and formats
kernel-reported matches.

The runner exposes this mechanism to agents through
\texttt{actplane run -- <cmd>}. It discovers the project policy file,
compiles and loads the policy before the target starts, seeds the target
process tree with the agent label, and records kernel matches in a
feedback file that compatible agent hooks can relay into the next turn.
This plumbing is intentionally thin: the policy decision remains in the
kernel, while the runner only connects project policy, target execution,
and agent-facing diagnostics.
